\documentclass[12pt]{article}

% Import preambles and macros for homework
\input{../../../../preambles/hw-preamble.tex}
\input{../../../../preambles/homework-macros.tex}
\input{../../../../preambles/homework-title.tex}
\input{../../../../preambles/homework-config.tex}

\renewcommand{\author}{Deepak Jassal}
\renewcommand{\authorlast}{Jassal}
\renewcommand{\coursename}{Advanced Linear Algebra}
\renewcommand{\coursecode}{MATH 326}
\renewcommand{\assignment}{Assignment 6}
\renewcommand{\instructor}{Dr. Edward Dobrowolski}
\renewcommand{\duedate}{March 27th\textsuperscript{th}, 2026}
\usepackage{tensor}

\begin{document}
\begin{titlepage}
	\centering
	\vspace*{2.0cm}	
	\pgfornament{84}\\
	{\LARGE \textsc{\coursename}\par}
	\vspace{0.5cm}
	{\large\coursecode\par}
    \vspace{0.5cm}
    {\large\instructor\par}
	\vspace{1.5cm}
	{\huge\bfseries\assignment\par}
	\vspace{1cm}  
	{\LARGE\itshape\author\par}
    \vspace{2cm}
	{\large\bfseries Due Date:\par}
	\vspace{0.5cm}
	{\Large \duedate}\\
	\pgfornament{84}
\end{titlepage}
\stepcounter{section}
\section*{Problem 1}
Let \(V = M_{3}(\mathbb{F})\), where \(\mathbb{F} = \mathbb{R}\) or \(\mathbb{C}\). State the standard basis of \(V\) and prove that it is a basis.
This problem appears because most answers on Midterm 2 about \(\dim(V)\) were wrong.
\begin{proof}
	The standard basis of $V$ is
\[
	\begin{aligned}
		\beta = \left\{
		\begin{bmatrix}
		1 & 0 & 0 \\
		0 & 0 & 0 \\
		0 & 0 & 0
		\end{bmatrix},
		\begin{bmatrix}
		0 & 1 & 0 \\
		0 & 0 & 0 \\
		0 & 0 & 0
		\end{bmatrix},
		\begin{bmatrix}
		0 & 0 & 1 \\
		0 & 0 & 0 \\
		0 & 0 & 0
		\end{bmatrix},\right. \\[4pt]
		\left.
		\begin{bmatrix}
		0 & 0 & 0 \\
		1 & 0 & 0 \\
		0 & 0 & 0
		\end{bmatrix},
		\begin{bmatrix}
		0 & 0 & 0 \\
		0 & 1 & 0 \\
		0 & 0 & 0
		\end{bmatrix},
		\begin{bmatrix}
		0 & 0 & 0 \\
		0 & 0 & 1 \\
		0 & 0 & 0
		\end{bmatrix},\right. \\[4pt]
		\left.
		\begin{bmatrix}
		0 & 0 & 0 \\
		0 & 0 & 0 \\
		1 & 0 & 0
		\end{bmatrix},
		\begin{bmatrix}
		0 & 0 & 0 \\
		0 & 0 & 0 \\
		0 & 1 & 0
		\end{bmatrix},
		\begin{bmatrix}
		0 & 0 & 0 \\
		0 & 0 & 0 \\
		0 & 0 & 1
		\end{bmatrix}
		\right\}
	\end{aligned}
\]
Each of these vectors is linearly independant (the sum of linear combinations of these 9 vectors can only be 0 if each coefficient is 0), the dimension of $M_3$ is 9. Since $\beta$ is a list of 9 linearly independant vectors in $M_3$ we have that $\beta$ is a basis (the standard basis) of $M_3$.
\end{proof}
\stepcounter{section}
\section*{Problem 2}
Let \(V = V[-1, 1]\) be a vector space over \(\mathbb{R}\) with inner product defined as 
\[
\langle f, g \rangle = \int_{-1}^{1} f(x)g(x) \, dx.
\]
By orthogonalizing the sequence \(\{1, x, x^{2}, x^{3}, \dots\}\) we obtain Legendre polynomials. The first five Legendre polynomials are 
\[
P_{0}(x) = 1,\quad P_{1}(x) = x,\quad P_{2}(x) = \frac{1}{2}(3x^{2} - 1),\quad P_{3}(x) = \frac{1}{2}(5x^{3} - 3x),\quad P_{4}(x) = \frac{1}{8}(35x^{4} - 30x^{2} + 3).
\]

\begin{enumerate}
    \item[(a)] Verify that \(\langle P_{2}, P_{3} \rangle = 0\).\\
    \textit{Solution.}
	\[
		\inner{P_2}{P_3}=\int_{-1}^{1} \left(\frac{1}{2}(3x^{2} - 1)\right)\left(\frac{1}{2}(5x^{3}-3x)\right) \, dx=\frac{1}{4}\int_{-1}^{1} 15x^5-14x^3+3x\,dx=0,
	\]
	since we are integrating an odd function over a symetric interval.
    \item[(b)] The listed above polynomials are standardized so that \(P_{n}(1) = 1\), however they are not orthonormal. Assuming their orthogonality, find the corresponding orthonormal polynomials.\\
    \textit{Solution.} To find the corresponding orthonormal polynomials we need to divide the polynomials by their norm.
	\[
		\norm{P_0}=\sqrt{\int_{-1}^{1}1\,dx}=\sqrt{2},\; \norm{P_1}=\sqrt{\int_{-1}^{1}x^2\,dx}=\sqrt{\frac{2}{3}},\;\norm{P_2}=\sqrt{\int_{-1}^{1}\left(\frac{1}{2}(3x^{2} - 1)\right)^2\,dx}=\sqrt{\frac{2}{5}}
	\]
	\[
		\norm{P_3}=\sqrt{\int_{-1}^{1}\left(\frac{1}{2}(5x^{3} - 3x)\right)^2\,dx}=\sqrt{\frac{2}{7}},\;\norm{P_4}=\sqrt{\int_{-1}^{1}\left(\frac{1}{8}(35x^{4} - 30x^{2} + 3)\right)^2\,dx}=\sqrt{\frac{2}{9}}.
	\]
	Thus, the corresponding orthonormal polynomials are 
	\[
		\hat{P_0}(x)=\frac{1}{\sqrt{2}},\;\hat{P_1}(x)=\left(\sqrt{\frac{3}{2}}\right)x,\;\hat{P_2}(x)=\left(\sqrt{\frac{5}{2}}\right)\frac{1}{2}(3x^{2} - 1),
	\]
	\[
		\hat{P_3}(x)=\left(\sqrt{\frac{7}{2}}\right)\frac{1}{2}(5x^{3} - x),\;\hat{P_4}(x)=\left(\sqrt{\frac{9}{2}}\right)\frac{1}{8}(36x^{4} - 30x^{2} + 3).
	\]
    \item[(c)] Let \(W = \operatorname{span}\{1, x, x^{2}, x^{3}, x^{4}\}\) and \(T: W \to \mathbb{R}\) be defined as \(T(f(x)) = f(1) + f'(2)\).\\
	\begin{enumerate}
		\item[(c1)] Verify that \(T\) is a linear functional on \(W\).\\
		\textit{Solution.}
		Let $f,g\in W$ and $a,b\in \R$ then, 
		\[
			T(af+bg)=af(1)+af'(2)+bg(1)+bg'(2)=aT(f)+bT(g).
		\]
		Thus, $T$ is a linear functional. 
		\item[(c2)] Find the Riesz vector of \(T\) using orthonormal basis found in (a).\\
		\textit{Solution.} The Riesz vector $v\in W$ is a vector such that $T(f)=\inner{f}{v}$ for all $f\in W$. We can use the orthonormal bassi we found above to $v$ in the following manner
		\[
			v=\sum_{i=0}^{4}\overline{T(\hat{P_i})}\hat{P_i}=\sum_{i=0}^{4}T(\hat{P_i})\hat{P_i}.
		\]
		Using a computer algebra system I computed $v$ to be
		\[
			v=\frac{1635}{8}-\frac{1215}{8}x-2100x^2+\frac{2065}{8}x^3+\frac{19845}{8}x^4.
		\]
		\newpage
		\item[(c3)] Verify that this vector works by checking it on \(f(x) = a + bx + cx^{2} + dx^{3} + ex^{4}\).\\
		\textit{Solution.}
		\[
			T(f)=a+2b+5c+13d+33e.
		\]
		Note that 
		\[
			f=\sum_{i=0}^{4}\inner{f}{\hat{P_i}}\hat{P_i}.
		\]
		We can compute the desired inner product as follows
		\[
			\inner{f}{v}=\sum_{i=0}^{4}T(\hat{P_i})\inner{f}{\hat{P_i}},
		\]
		by linearity of $T$ and the above formula for $f$ we have
		\[
			\inner{f}{v}=T\left(\sum_{i=0}^{4}\inner{f}{\hat{P_i}}\hat{P_i}\right)=T(f).
		\]
	\end{enumerate}
\end{enumerate}
You can use software for the calculations to save time.

\stepcounter{section}
\section*{Problem 3}
Use Gram-Schmidt procedure to find the orthonormal basis of \(\mathbb{R}^{3}\) by orthogonalizing the basis 
\[
\{\langle 1, 1, 1 \rangle, \langle 1, 1, 0 \rangle, \langle 1, 0, 0 \rangle\}.
\]
Note: \(\{e_{1}, e_{2}, e_{3}\}\) is not the solution.\\
\textit{Solution.}
\[
	u_1=v_1=\langle 1,1,1\rangle,\;\norm{u_1}=\sqrt{3},\,e_1'=\frac{1}{\sqrt{3}}\vect{1,1,1}.
\]
\begin{align*}
	u_2&=v_2-\frac{\inner{v_2}{u_1}}{\norm{u_1}^2}u_1\\
	&=\vect{1,1,0}-\frac{2}{3}\vect{1,1,1}\\
	&=\vect{\frac{1}{3},\frac{1}{3},-\frac{2}{3}}\\
	\norm{u_2}&=\sqrt{\frac{6}{9}}\\
	e_2&=\frac{1}{\sqrt{6}}\vect{1,1,-2}
\end{align*}
\begin{align*}
	u_3&=v_3-\frac{\inner{v_3}{u_1}}{\norm{u_1}^2}u_1-\frac{\inner{v_3}{u_2}}{\norm{u_2}^2}u_2\\
	&=\vect{1,0,0}-\frac{1}{3}\vect{1,1,1}-\frac{1}{6}\vect{1,1,-2}\\
	&=\vect{\frac{1}{2},-\frac{1}{2},0}\\
	\norm{u_3}&=\frac{1}{\sqrt{2}}\\
	e_3&=\frac{1}{\sqrt{2}}\vect{1,-1,0}.
\end{align*}
So the orthonormal basis is 
\[
	\frac{1}{\sqrt{3}}\vect{1,1,1},\;\frac{1}{\sqrt{6}}\vect{1,1,-2},\;\frac{1}{\sqrt{2}}\vect{1,-1,0}.
\]

\stepcounter{section}
\section*{Problem 4}
Let \(U = \operatorname{span}\{\langle 1, 1, 1 \rangle, \langle 1, 1, 0 \rangle\}\). (Note: These are the first two vectors from Problem 3). Let \(P_{U}: \mathbb{R}^{3} \to \mathbb{R}^{3}\) be the projection on the plane \(U\).

\begin{enumerate}[label=(\alph*)]
    \item Find the matrix of \(P_{U}\).\\
    \textit{Solution.} We can use the orthonormal basis found above to calculate $P_U$.
	\begin{align*}
		P_U&=e_1e_1^T+e_2e_2^T\\
		&=\frac{1}{3}\vect{1,1,1}^T\vect{1,1,1}+\frac{1}{6}\vect{1,1,-2}^T\vect{1,1,-2}\\
		&=\frac{1}{3}
		\begin{bmatrix}
			1&1&1\\
			1&1&1\\
			1&1&1\\
		\end{bmatrix}
		+\frac{1}{6}
		\begin{bmatrix}
			1&1&-2\\
			1&1&-2\\
			-2&-2&4
		\end{bmatrix}\\
		&=\frac{1}{2}
		\begin{bmatrix}
			1&1&0\\
			1&1&0\\
			0&0&2
		\end{bmatrix}.
	\end{align*}
    \item Describe \(U^{\perp}\) geometrically and find the matrix of \(P_{U^{\perp}}\).\\
    \textit{Solution.} $U^\perp$ is spanned by $u_3'$, that is $\vect{1,-1,0}$. We can calculate the matrix for with the following formula
	\begin{align*}
		P_{U^\perp}&=I-P_U\\
		&=\frac{1}{2}
		\begin{bmatrix}
			1&-1&0\\
			-1&1&0\\
			0&0&0
		\end{bmatrix}
	\end{align*}
	\newpage
    \item We know that \(\mathbb{R}^{3} = U \oplus U^{\perp}\). Let \(\mathbf{v} = \langle 1, 2, 3 \rangle\). Find \(\mathbf{u} \in U\) and \(\mathbf{w} \in U^{\perp}\) such that \(\mathbf{v} = \mathbf{u} + \mathbf{w}\).\\
    \textit{Solution.}
	\[
		\mathbf{u}=P_U\mathbf{v}=
		\frac{1}{2}
		\begin{bmatrix}
			1&1&0\\
			1&1&0\\
			0&0&2
		\end{bmatrix}
		\begin{bmatrix}
			1\\2\\3
		\end{bmatrix}
		=
		\begin{bmatrix}
			\frac{3}{2}\\\frac{3}{2}\\3
		\end{bmatrix}
	\]
	and
	\[
		\mathbf{w}=P_{U^\perp}\mathbf{v}=
		\frac{1}{2}
		\begin{bmatrix}
			1&-1&0\\
			-1&1&0\\
			0&0&0
		\end{bmatrix}
		\begin{bmatrix}
			1\\2\\3
		\end{bmatrix}=
		\begin{bmatrix}
			-\frac{1}{2},\frac{1}{2},0
		\end{bmatrix}.
	\]
\end{enumerate}



\end{document}